\documentclass{situationist-preprint}

\title{The Ship of Theseus Problem Settled}
\msid{MET-07-a}
\msfamily{MET-07}
\msdiscipline{Metaphysics}
\mstype{Principal manuscript}
\mskeywords{Ship of Theseus; artefact identity; persistence; maintenance; only-$x$-and-$y$ principle; closest continuer}

\begin{document}
\maketitle

\begin{abstract}
A ship has its planks replaced one by one; the discarded planks are stored and later reassembled in their original arrangement. Which of the two resulting ships is the ship we started with? We argue that the continuously repaired ship is the original, and that the reassembled ship is a distinct ship composed of the original planks. The argument rests on three principles about artefacts: that they persist through the replacement of parts when that replacement is a stage in their maintenance; that their identity over time is fixed by the maintenance relation rather than by material continuity alone; and that this relation is intrinsic to the processes that connect earlier and later ships. The resulting account respects the only-$x$-and-$y$ principle, which closest-continuer theories violate, and explains why a disassembled watch can survive its repair while Hobbes's reassembled ship does not survive its reconstruction.
\end{abstract}

\section{The puzzle}

Plutarch reports that the ship in which Theseus returned from Crete was preserved by the Athenians, who \enquote{took away the old timbers from time to time, and put new and sound ones in their place}, so that the ship became, for philosophers, a standing example of the logical question of growth \citep{Plutarch}. Hobbes added the twist that turns the example into a genuine puzzle \citep{Hobbes}. Suppose that someone keeps every discarded plank and, once all of them have been removed, reassembles them in their original order. There are now two ships, and each has a credible claim to be the ship of Theseus. The repaired ship has the better claim from the point of view of continuity of form and location; the reassembled ship has the better claim from the point of view of material composition.

Let $S_0$ be the ship at the start, $S_R$ the repaired ship at the end, and $S_A$ the reassembled ship. The positions in the literature can be sorted according to how they answer the question whether $S_0 = S_R$ or $S_0 = S_A$.
\begin{enumerate}[label=(\roman*)]
\item \emph{Material views} identify $S_0$ with $S_A$, on the ground that a ship is, or is essentially tied to, the planks that compose it. In its strongest form this is Chisholm's mereological essentialism, on which no composite object can survive the loss of any of its parts \citep{Chisholm}. On that view $S_0$ does not survive even the first replacement, and the question is not so much answered as dissolved.
\item \emph{Continuity views} identify $S_0$ with $S_R$, on the ground that the persistence of a ship is a matter of spatio-temporal and causal continuity under the sortal \enquote{ship} \citep{Wiggins1980,Wiggins2001}.
\item \emph{Closest-continuer views} hold that $S_0$ is whichever successor stands in the closest relation to it, provided that relation is close enough \citep{Nozick}. Whether $S_R$ or $S_A$ is the original then depends on a comparative weighing of continuity and composition.
\item \emph{Deflationary views} hold that there is no fact of the matter, or that the question is verbal, or that it concerns only which of two overlapping four-dimensional objects deserves the name \citep{Lewis1976,Sider2001}.
\end{enumerate}

We defend a version of the continuity view. The standard objection to it is that it seems unable to explain a case that looks structurally identical: a watch taken apart for cleaning, its parts kept in a tray, and then reassembled. Most people judge that the reassembled watch is the same watch, which seems to support the material view. A continuity theorist must either deny this judgement or explain why the watch and the ship differ. Explaining the difference without making the identity of the ship depend on extrinsic facts is the main task of this manuscript.

\section{Framework}

We work with ordinary objects and with three-dimensional persistence. Nothing in the argument depends on this choice; Section~6 shows how the result transfers to a four-dimensionalist framework.

\begin{definition}
An \emph{artefact of kind $K$} is an object that exists in virtue of being made, used or maintained as a $K$ by agents who have the concept of $K$.
\end{definition}

The definition follows the broadly intention-dependent conception of artefacts defended by \citet{Baker2007} and \citet{Thomasson2007}. It does not say that artefacts are mind-dependent in any problematic sense: a ship is made of wood and floats whether or not anyone thinks about it. It says that what makes some planks a ship, rather than a heap of wood that happens to float, involves a practice of building and using ships.

\begin{definition}
A process $\pi$ connecting an object $x$ at $t$ to an object $y$ at $t'$ is a \emph{maintenance process for $x$} if and only if every stage of $\pi$ is carried out with the aim of keeping $x$ in service as a $K$, and at every stage the object undergoing the process retains the organization in virtue of which it functions as a $K$, or is restored to that organization within the course of a single repair.
\end{definition}

The last clause is what accommodates the watch. A watch in pieces in the repairer's tray does not, at that moment, have the organization of a working watch; but it is being restored to that organization within a single repair, and the process is aimed at keeping \emph{that} watch in service. Whether the watch exists while its parts lie in the tray is a question for the theory of composition. On the account of MET-06 the parts compose nothing at that moment, and the watch exists intermittently; nothing in the present argument requires that it exist continuously.

We now state the three principles on which the argument depends.

\begin{principle}[Tolerance]\label{p:tolerance}
An artefact of kind $K$ can survive the replacement of any of its parts, and of all its parts over time, provided each replacement is a stage of a maintenance process for that artefact.
\end{principle}

\begin{principle}[Maintenance]\label{p:maintenance}
If $x$ at $t$ is an artefact of kind $K$, then $y$ at $t'$ is identical to $x$ if and only if $y$ is a $K$ at $t'$ and $x$ at $t$ is connected to $y$ at $t'$ by a maintenance process for $x$.
\end{principle}

\begin{principle}[Intrinsicness]\label{p:intrinsic}
Whether $x$ at $t$ is identical to $y$ at $t'$ depends only on $x$, $y$ and the processes that connect them.
\end{principle}

Principle~\ref{p:intrinsic} is a version of the only-$x$-and-$y$ principle \citep{Wiggins1980,Noonan1985}. It rules out theories on which the identity of $x$ and $y$ depends on the existence of a third object, as in closest-continuer theories. The principle is widely accepted for the identity of material objects, and it is the principle that the obvious solutions to the ship puzzle tend to violate.

\section{The argument}

The case has two branches. In the first, a crew replaces the planks of $S_0$ one at a time, over many years, in order to keep it seaworthy. In the second, a collector gathers each discarded plank as it is removed, stores the planks, and eventually assembles them according to the original plan.

\begin{lemma}\label{l:repair}
The process connecting $S_0$ to $S_R$ is a maintenance process for $S_0$.
\end{lemma}

\begin{proof}
Each stage of the process is the replacement of a plank by a crew whose aim is to keep the ship in service. At every stage, the object being repaired has the organization of a seaworthy ship, since only one plank is missing at a time and it is immediately replaced. Both clauses of Definition~2.2 are therefore satisfied.
\end{proof}

\begin{lemma}\label{l:assembly}
The process connecting $S_0$ to $S_A$ is not a maintenance process for $S_0$.
\end{lemma}

\begin{proof}
Consider the collector's activity at any stage before the final assembly. Its aim is not to keep $S_0$ in service: the collector neither sails nor repairs $S_0$, and the planks in storage are not being restored to the organization of a ship within the course of a single repair, since no repair of any ship is in progress. The first clause of Definition~2.2 fails. The second fails as well: for the whole period of storage, the planks have no ship-organization and are not in the course of regaining it.
\end{proof}

It is important that Lemma~\ref{l:assembly} does not appeal to the existence of $S_R$. The collector's process fails to be a maintenance process for $S_0$ because of what the collector does and aims at, not because someone else happens to be repairing the ship in the harbour. We return to this point in the next section.

\begin{theorem}\label{t:main}
$S_0 = S_R$ and $S_0 \neq S_A$.
\end{theorem}

\begin{proof}
By Lemma~\ref{l:repair} and Principle~\ref{p:maintenance}, $S_0 = S_R$, since $S_R$ is a ship and is connected to $S_0$ by a maintenance process for $S_0$. By Lemma~\ref{l:assembly} and the right-to-left direction of Principle~\ref{p:maintenance}, $S_A$ is not identical to $S_0$, since no maintenance process for $S_0$ connects them. Principle~\ref{p:tolerance} guarantees that the conclusion $S_0 = S_R$ is not blocked by the fact that $S_R$ shares no parts with $S_0$.
\end{proof}

\begin{corollary}\label{c:new}
$S_A$ is a ship that began to exist at the time of its assembly, and is composed of the planks that once composed $S_0$.
\end{corollary}

\begin{proof}
$S_A$ is a ship at the end of the process. By Theorem~\ref{t:main} it is not $S_0$, and the only ship whose maintenance history could have reached back before the assembly is $S_0$. Hence $S_A$ has no history as a ship before its assembly.
\end{proof}

\section{Why the account respects intrinsicness}

Closest-continuer theories give the right verdict in Hobbes's case, but for the wrong reason. On such a theory, $S_A$ would have been $S_0$ if no repaired ship had existed; it fails to be $S_0$ only because $S_R$ is a closer continuer. Identity then depends on a third object, in violation of Principle~\ref{p:intrinsic}. This is the familiar objection to closest-continuer theories \citep{Noonan1985}, and it applies with full force to artefacts.

Our account avoids it. Consider the variant in which the crew does not repair $S_0$ but simply lets it rot, while the collector gathers its falling planks and later reassembles them. Is $S_A$ now identical to $S_0$? Lemma~\ref{l:assembly} applies unchanged, since its proof makes no reference to the repaired ship. So $S_A$ is still not $S_0$: the original ship ceased to exist when it decayed beyond the point at which a single repair could restore it. This verdict may seem harsh, but it is exactly what the only-$x$-and-$y$ principle requires. The identity of $S_0$ and $S_A$ cannot depend on whether a crew happened to be at work on another ship.

Contrast the watch. When a repairer takes a watch apart, the process connecting the watch before and after is aimed at keeping that watch in service, and the parts in the tray are being restored to the organization of a watch within the course of a single repair. Definition~2.2 is satisfied, and Principle~\ref{p:maintenance} delivers identity. The watch and the ship differ, not because one has a rival and the other does not, but because one undergoes a repair and the other undergoes a reconstruction. The difference lies entirely in the processes that connect the earlier object with the later one, which is where Principle~\ref{p:intrinsic} says it should lie.

\section{Objections}

\begin{objection}
The account makes the identity of a ship depend on the aims of its repairers. But identity is a matter of what the world is like, not of what people intend.
\end{objection}

\begin{reply}
The account makes the identity of a ship depend on facts about maintenance, and these include facts about intentions. This is a consequence of a plausible view about what artefacts are (Definition~2.1), not an additional assumption. If being a ship involves a practice of making and using ships, it is unsurprising that the persistence conditions of ships also involve that practice. Natural objects are a different matter, and nothing said here applies to them.
\end{reply}

\begin{objection}
The material view has a strong intuitive basis. If the collector shows us $S_A$ and says that every plank in it was in Theseus's ship, arranged exactly as it was, it is hard to deny that this is Theseus's ship.
\end{objection}

\begin{reply}
The intuition is real, but it tracks a different property. $S_A$ is made of the very wood of Theseus's ship, and has historical and perhaps commercial value for that reason. What the intuition supports is that $S_A$ is \emph{composed of} what once composed $S_0$, which Corollary~\ref{c:new} affirms. It does not support identity. The same distinction applies to relics: a reliquary may contain all the bones of a saint without being the saint.
\end{reply}

\begin{objection}
Principle~\ref{p:tolerance} leads to a sorites. If a ship survives the replacement of one plank, it survives the replacement of two, and so on. But surely a ship that has had all its planks replaced at once, by a crew that tears it apart and builds another one beside it, is not the same ship.
\end{objection}

\begin{reply}
We agree with the verdict, and the account delivers it. Tearing a ship apart and building another one is not a maintenance process, since the object undergoing the process loses its ship-organization and is not restored to it within the course of a single repair. What Principle~\ref{p:tolerance} licenses is gradual replacement within maintenance, and the sorites does not get started because the condition on each step is not simply that one more plank has been replaced.
\end{reply}

\begin{objection}
The question is verbal. There are two candidate referents for \enquote{the ship of Theseus}, and we can use the expression for either.
\end{objection}

\begin{reply}
We can, but the practices that give the expression its meaning settle which use is correct. Ships are named, registered, owned, insured and sailed; all these practices follow the maintained ship. A port authority that recorded $S_A$ as the continuation of $S_0$ would be making an error, not choosing a convention. Companion manuscript MET-07-b examines the relation between ownership and identity in detail.
\end{reply}

\section{Four-dimensionalism}

On a perdurantist view, $S_0$, $S_R$ and $S_A$ are temporal parts of four-dimensional objects, and the question becomes which of the overlapping worms that contain $S_0$ as a stage is the referent of \enquote{the ship of Theseus} \citep{Lewis1976,Sider2001}. Theorem~\ref{t:main} can be restated in these terms: the worm whose stages are unified by the maintenance relation is a ship; the worm that runs from $S_0$ through the stored planks to $S_A$ is not a ship, since some of its stages are heaps of planks that are not in the course of being restored to ship-organization. The perdurantist who accepts Definition~2.1 therefore reaches the same verdict, and the result does not depend on the choice between three-dimensionalism and four-dimensionalism.

\section{Conclusion}

The continuously repaired ship is the ship of Theseus. The reassembled ship is a different ship, which began to exist when it was assembled and is composed of the original planks. The verdict follows from three principles about artefacts, one of which is the only-$x$-and-$y$ principle; it explains the difference between Hobbes's ship and a repaired watch; and it is available to three-dimensionalists and four-dimensionalists alike. The case that has served for two millennia as an emblem of unanswerable questions about identity turns out to have an answer, once we attend to what ships are.

\begin{thebibliography}{99}\small
\bibitem[Baker(2007)]{Baker2007} Baker, L.~R. (2007). \emph{The Metaphysics of Everyday Life: An Essay in Practical Realism}. Cambridge University Press.
\bibitem[Chisholm(1976)]{Chisholm} Chisholm, R.~M. (1976). \emph{Person and Object: A Metaphysical Study}. Open Court.
\bibitem[Hobbes(1655)]{Hobbes} Hobbes, T. (1655). \emph{De Corpore}, Part II, Chapter 11, §7.
\bibitem[Lewis(1976)]{Lewis1976} Lewis, D. (1976). Survival and identity. In A.~O. Rorty (Ed.), \emph{The Identities of Persons} (pp.~17--40). University of California Press.
\bibitem[Noonan(1985)]{Noonan1985} Noonan, H.~W. (1985). The only $x$ and $y$ principle. \emph{Analysis}, 45(2), 79--83.
\bibitem[Nozick(1981)]{Nozick} Nozick, R. (1981). \emph{Philosophical Explanations}. Harvard University Press.
\bibitem[Plutarch(n.d.)]{Plutarch} Plutarch. \emph{Life of Theseus}, 23.1.
\bibitem[Sider(2001)]{Sider2001} Sider, T. (2001). \emph{Four-Dimensionalism: An Ontology of Persistence and Time}. Oxford University Press.
\bibitem[Thomasson(2007)]{Thomasson2007} Thomasson, A.~L. (2007). \emph{Ordinary Objects}. Oxford University Press.
\bibitem[Wiggins(1980)]{Wiggins1980} Wiggins, D. (1980). \emph{Sameness and Substance}. Blackwell.
\bibitem[Wiggins(2001)]{Wiggins2001} Wiggins, D. (2001). \emph{Sameness and Substance Renewed}. Cambridge University Press.
\end{thebibliography}

\end{document}
